\section{Involution entre les coordonnées de résidu et de quotient}
\label{sec:dualidad-visible-memoria}

La division euclidienne utilisée dans \(\mathcal A_9^\#\) s'écrit, pour un entier dans
la carte nonadique,

\[
n=9Q+r,
\qquad r\in\{1,\ldots,9\}.
\]

La coordonnée \(r\) est le représentant observable et \(Q\) enregistre le
nombre de franchissements de frontière. Pour formuler une involution qui échange
les deux coordonnées, il est nécessaire de passer à une extension stable sous
l'échange. Notons \(\bar r\in\{0,\ldots,8\}\) la classe résiduelle,
avec \(\bar r=0\) lorsque le représentant est \(r=9\), et considérons

\[
\mathscr D_{\mathrm{dual}}=\mathbb Z^2\times\mathbb R_{>0}.
\]

Les deux entiers de \(\mathscr D_{\mathrm{dual}}\) ne sont pas restreints à la section canonique des
représentants ; cet élargissement est nécessaire puisqu'un quotient \(Q\) peut
se trouver hors de l'intervalle \(0,\ldots,8\) lorsqu'il occupe la première coordonnée.

\begin{definition}[Forme d'énergie de rayon]
Pour \(R_0>0\) et \((m,w)\in\mathbb Z^2\), nous définissons
\[
E_{R_0}(m,w)=\frac{m^2}{R_0^2}+w^2R_0^2.
\]
La transformation de dualité est
\[
\mathsf T(m,w,R_0)=(w,m,R_0^{-1}).
\]
\end{definition}

\begin{theorem}[Dualité résidu--quotient]
\label{thm:dualidad-visible-memoria}
L'application \(\mathsf T\) est une involution de \(\mathscr D_{\mathrm{dual}}\) et préserve
la forme d'énergie au sens
\[
E_{R_0}(m,w)=E_{R_0^{-1}}(w,m).
\]
\end{theorem}

\begin{proof}
La seconde application de \(\mathsf T\) échange de nouveau \(m,w\) et
inverse une seconde fois \(R_0\), de sorte que
\[
\mathsf T^2(m,w,R_0)=(m,w,R_0).
\]
De plus,
\[
E_{R_0^{-1}}(w,m)
=\frac{w^2}{R_0^{-2}}+m^2R_0^{-2}
=w^2R_0^2+\frac{m^2}{R_0^2}
=E_{R_0}(m,w).\qedhere
\]
\end{proof}

L'insertion d'un état nonadique dans le domaine élargi est

\[
(r,Q)\longmapsto(\bar r,Q,R_0).
\]

Le transformé \((Q,\bar r,R_0^{-1})\) n'appartient pas nécessairement à la
section \(0\leq\bar r\leq8\). Pour cette raison, \(\mathsf T\) est une
involution exacte du domaine élargi, non un endomorphisme de l'atlas résiduel
à 81 positions. Le point fixe de la transformation du paramètre est
\(R_0=1\) ; en ce point, l'échange des deux coordonnées préserve la même forme
quadratique.

L'identité admet une réalisation opératorielle. Soit
\(\mathcal H_{\rm lat}=\ell^2(\mathbb Z^2)\), de base orthonormée
\(\{|m,w\rangle\}\). Définissons l'opérateur diagonal

\[
H(R_0)|m,w\rangle=E_{R_0}(m,w)|m,w\rangle
\]

sur son domaine naturel, et l'opérateur unitaire

\[
U_T|m,w\rangle=|w,m\rangle.
\]

\begin{corollary}[Équivalence unitaire des rayons réciproques]
\label{cor:dualidad-unitaria}
Pour tout \(R_0>0\),
\[
U_TH(R_0)U_T^{-1}=H(R_0^{-1}).
\]
\end{corollary}

\begin{proof}
Sur un vecteur de la base,
\[
U_TH(R_0)U_T^{-1}|m,w\rangle
=E_{R_0}(w,m)|m,w\rangle.
\]
Le \cref{thm:dualidad-visible-memoria} donne
\(E_{R_0}(w,m)=E_{R_0^{-1}}(m,w)\), qui est l'action de
\(H(R_0^{-1})\).
\end{proof}

Cette conjugaison a la forme mathématique de la dualité de rayon dans une
compactification circulaire, dans des unités où l'échelle de corde a été
normalisée à un \cite{Giveon1994}. L'identification physique supplémentaire
attribue \(m\) au nombre d'impulsion, \(w\) au nombre d'enroulement et \(R_0\)
au rayon de compactification. Le théorème n'introduit ni oscillateurs, ni champs,
ni supersymétrie, ni tension de corde ; il démontre l'involution de réseau et
l'équivalence unitaire des opérateurs diagonaux déclarés.
